\documentclass[a4paper,10pt]{article}
\usepackage[utf8]{inputenc}


\usepackage{lipsum}
\usepackage[margin=25mm]{geometry}
\usepackage{mdframed}
\usepackage{xcolor}
\usepackage{hyperref}

\usepackage[
AsteriskCount=5,
BeforeSpace=0pt,
AfterSpace=0pt,
Symbol=*
]{ornamental-break-asterisks}


%opening
\title{ornamental-break-asterisks package}
\author{Fernando Pujaico Rivera}



\begin{document}

\maketitle

\begin{abstract}
Creates macros to show ornamental breaks.
\end{abstract}

%%%%%%%%%%%%%%%%%%%%%%%%%%%%%%%%%%%%%%%%%%%%%%%%%%%%%%%%%%%%%%%%%%%%%%%%%%%%%%%%%%%%%%%%%%%%%%%%%
\section{Install package}
Put the ornamental-break-asterisks.sty file in any of these locations
\begin{itemize}
\item Put the \texttt{ornamental-break-asterisks.sty} file in the same path of main tex file, or.
\item Execute the commmand:
\begin{mdframed}[backgroundcolor=black!10,rightline=false,leftline=false]
\begin{verbatim}
kpsewhich -var-value=TEXMFHOME
\end{verbatim}
\end{mdframed}
and this returns the path of local tex files. By example, if returns 
\begin{mdframed}[backgroundcolor=black!10,rightline=false,leftline=false]
\begin{verbatim}
/home/username/texmf
\end{verbatim}
\end{mdframed}
then, put the \texttt{ornamental-break-asterisks.sty} file in the directory.
\begin{mdframed}[backgroundcolor=black!10,rightline=false,leftline=false]
\small
\begin{verbatim}
/home/username/texmf/tex/latex/ornamental-break-asterisks/ornamental-break-asterisks.sty
\end{verbatim}
\end{mdframed}
\end{itemize}



%%%%%%%%%%%%%%%%%%%%%%%%%%%%%%%%%%%%%%%%%%%%%%%%%%%%%%%%%%%%%%%%%%%%%%%%%%%%%%%%%%%%%%%%%%%%%%%%%

\section{Load the package}
To load the package use the next command in the preamble of main tex document.

\begin{mdframed}[backgroundcolor=black!10,rightline=false,leftline=false]
\begin{verbatim}
\usepackage[
AsteriskCount=5,
BeforeSpace=0pt,
AfterSpace=0pt,
Symbol=*
]{ornamental-break-asterisks}
\end{verbatim}
\end{mdframed}


%%%%%%%%%%%%%%%%%%%%%%%%%%%%%%%%%%%%%%%%%%%%%%%%%%%%%%%%%%%%%%%%%%%%%%%%%%%%%%%%%%%%%%%%%%%%%%%%%

\section{Using the package}
Next, I will provide examples that demonstrate the use of all the macros included in the package. 
This will help illustrate how each macro functions within the context of the package and provide 
practical insights into their applications.

%%%%%%%%%%%%%%%%%%%%%%%%%%%%%%%%%%%%%%%%%%%%%%%%%%%%%%%%%%%%%%%%%%%%%%%%%%%%%%%%%%%%%%%%%%%%%%%%%
%%%%%%%%%%%%%%%%%%%%%%%%%%%%%%%%%%%%%%%%%%%%%%%%%%%%%%%%%%%%%%%%%%%%%%%%%%%%%%%%%%%%%%%%%%%%%%%%%

\begin{verbatim}
\lipsum[1][1-4]

\OrnamentalBreak

\lipsum[1][1-4]
\end{verbatim}

\lipsum[1][1-4]

\OrnamentalBreak

\lipsum[1][1-4]


\end{document}
